\documentclass{article}
\usepackage{amsmath,amssymb}
\usepackage[vocab={recurrences}]{leanify}
\newtheorem{lemma}{Lemma}
\begin{document}

\section{Notation}

We write $A_n(a,d)$ for the $n$th term of the arithmetic progression with first term $a$ and difference $d$.
We write $S_n(a,d)$ for the sum of the first $n$ terms of the arithmetic progression with first term $a$ and difference $d$.
We write $G_n(a,q)$ for the $n$th term of the geometric progression with first term $a$ and ratio $q$.
We write $T_n(a,q)$ for the sum of the first $n$ terms of the geometric progression with first term $a$ and ratio $q$.
We write $\lambda(p,q)$ for the larger characteristic root of the recurrence with coefficients $p$ and $q$.
For every real number $n$, let $u(n)\coloneqq 3n+1$.

\section{Progressions}

\begin{lemma}\label{arithmetic}
$u$ is an arithmetic progression with difference $3$, and $A_4(1,3)=13$.
\end{lemma}

\begin{lemma}\label{sum}
Let $n$ be a nonnegative integer. Then $S_n(1,2)=n^2$.
\end{lemma}

\begin{lemma}\label{geometric}
Let $a:\mathbb{R}\to\mathbb{R}$ be a function and let $q$ be a real number. If $a$ is a geometric progression with ratio $q$, then $a(2)=q^2a(0)$, and $G_2(1,q)=q^2$.
\end{lemma}

\section{Recurrences of order two}

\begin{lemma}\label{fibonacci}
Let $a:\mathbb{R}\to\mathbb{R}$ be a function. If $a$ is a solution of the Fibonacci recurrence, then $a$ is a solution of the linear recurrence with coefficients $1$ and $1$.
\end{lemma}

\begin{lemma}\label{root}
$\lambda(1,1)=\frac{1+\sqrt5}{2}$, and $\lambda(1,1)^2=\lambda(1,1)+1$.
\end{lemma}

\section{Closed forms}

\begin{lemma}\label{closed}
Let $n$ be a nonnegative integer and let $a$, $d$ and $q$ be real numbers with $q\neq1$. Then $S_n(a,d)=na+\frac{n(n-1)d}{2}$ and $T_n(a,q)=\frac{a(1-q^n)}{1-q}$.
\end{lemma}

\end{document}
